La distribución de los rendimientos diarios evidencia diferencias importantes en el nivel de riesgo entre las empresas con mayor rendimiento acumulado. BNTX presenta la mayor dispersión de rendimientos, con un rango intercuartílico aproximado entre -2.27% y 2.57%, lo que indica que sus variaciones diarias habituales fueron significativamente superiores al resto de compañías analizadas. Asimismo, concentra los valores atípicos más extremos, alcanzando aproximadamente pérdidas cercanas al -35% y ganancias superiores al 70% en un solo día, reflejando episodios de elevada incertidumbre y fuertes reacciones del mercado.
PDD registra la segunda mayor dispersión, con un rango intercuartílico aproximado entre -2.11% y 2.29%, acompañado de valores atípicos que alcanzan alrededor del -25% y 60%. Este comportamiento confirma que ambas compañías estuvieron expuestas a fluctuaciones diarias considerablemente superiores al promedio, por lo que, aunque ofrecieron elevados rendimientos acumulados, también implicaron un mayor nivel de riesgo para los inversionistas.
En contraste, NVIDIA, a pesar de haber obtenido el mayor rendimiento acumulado del período, presenta una distribución diaria relativamente más compacta, con un rango intercuartílico aproximado entre -1.54% y 1.99%. Esto sugiere que su crecimiento no dependió únicamente de movimientos extremos, sino de una apreciación más consistente a lo largo del tiempo. De forma similar, FICO exhibe una de las distribuciones más estrechas del grupo, indicando una menor volatilidad diaria y un comportamiento relativamente más estable, aun cuando también registró rendimientos acumulados sobresalientes.